\section{Introduction}

% Providing highly reliable while quickly generalizing new tasks is important in automation. Recent learning techniques, LLM increase capability, VLM increase perception, VLA is learned. GoFE is alreayd great.geeralization sometimes comes with the cost of reliabiltiy and specialization. Goal is to harness the existing capabilities of generalizaton without sacrificacing reliabiltiy of industrial automation. 
% Learning to combine different methods have made possible. 


A majority of robot learning research focuses on generalist robotics, where a robot must perform a broad variety of tasks.  Most of this research has focused on model-free end-to-end Vision-Language-Action (VLA) models ~\cite{kim24openvla, pi0_2024, open_x_embodiment_rt_x_2023, fang2026molmoact2, black2025pi_, zawalski2025robotic}.

However fully generalist robots struggle to achieve commercial / industrial levels of reliability \cite{gao2025taxonomy} and there is increasing interest in how robot learning could be useful for more specialized task classes~\cite{goldberg2026generalists}. In this paper, we identify a class of ``Variational Automation (VA)'' tasks that differ from ``fixed automation'' (which blindly repeats the same motions e.g., spot welding or spray painting).  In VA tasks, a robot persistently performs varying instances of a task with non-trivial variation in the geometry and pose of objects (e.g., to sort packages, make coffee in a cafe, or build sandwiches in a commerical kitchen).
  
Today, fixed automation is set up and tuned by humans using rigorous, classical engineering methods in logistics, service, agriculture and manufacturing \cite{solowjow_indust_2020, adebola2024automating, xie_energy_2025, adebolagarden2023}. This provides high reliability and throughput, and human effort can be justified and amortized over years of repetitive performance. But even more human effort is required to set up and tune VA tasks for reliable performance.

Recent advances in Large Language Models (LLM)~\cite{team2026qwen3, yang2025qwen3,guo2025deepseek} and Vision-Language Models (VLM)~\cite{radford2021learning, bordes2024introduction, geminiroboticsteam2025geminiroboticsbringingai, huang2023voxposer,ning2025prompting} are rapidly improving agentic coding~\cite{claudecode}, semantic reasoning~\cite{chenrobo2vlm} and zero-shot generalization~\cite{ji2025robobrain}.

Agentic coding for robotics may have the potential to integrate model-based and model-free paradigms to rapidly generate interpretable and reliable robot control systems. Agentic coding is new and advancing rapidly, but remains prone to hallucinations and constraint violations. So a primary challenge is to develop effective methods for ``steering'' coding agents to produce desired outputs.  Coding agents are controlled with a ``harness'' -- a body of text provided as a prompt to the agents that describe available software resources, constraints, and performance metrics~\cite{anthropicharness}.  The harness also provides interfaces that allow the LLM to compile, execute, and manage generate-execute-observe-evaluate loops.

The use of LLM coding agents for robotics was explored as early as 2022 with a series of papers on 
Code-as-Policy (CaP) ~\cite{codeaspolicies2022, YinG-RSS-25}.
General coding agents have improved dramatically since then by training on increasingly large coding datasets.  The recent CaP-X paper provides an open suite of CaP benchmarks, agents, and encouraging results using coding agents that were state-of-the-art as of January 2026 ~\cite{fu2026cap}.

These single-agent CaP approaches are somewhat unstructured, prompting the coding agent to generate free-form python code.  For more complex tasks, the context window can grow sufficiently large so that it is difficult for the agent to obey constraints and converge on reliable code.  Individual agents are also prone to hallucinations and ``cheating'', where they make up nonexistent skills or create trivial success metrics to artificially ''solve'' robot coding tasks.  These challenges are exacerbated with \emph{multi-agent} coding systems, where multiple agents operate in parallel.

To address these issues, we propose Graph-as-Policy (GaP), a harness structure that encapsulates modular robot functions into independent ``nodes'' in a directed graph that can be managed by a hierarchical multi-agent system where optimizing each node can be assigned to a specific agent.  This structure also helps limit the size of each agent's context window and separates the generation of graph elements from the testing of graph elements to reduce incentives for individual coding agents to ``cheat'' in order to achieve requested objectives.

GaP is inspired in part by the architecture of robot Task and Motion Planning (TAMP) systems~\cite{kaelbling2011hierarchical,garrett2021integrated, shen2026tiptop} which use graphical hierarchies to ensure safety and the Robot Operating System (ROS)\cite{macenski2022robot} which is based on a computation graph structure.  Graphical models support modularity, reuse, and composability of functions in non-robotic and robotic applications.   In GaP, a robot policy is a computation graph composed of atomic ``skill'' nodes, such as retrieving a camera frame, running inference on a perception model, or planning and executing a motion trajectory. 

Given a VA task description, GaP starts with an Orchestration Agent that partitions a VA task into functional segments and instructs appropriate Skill Agents to synthesize localized, functional subgraphs of atomic nodes for its assigned segment. The orchestrator then aggregates and wires these subgraphs together into an executable computation graph.  GaP then initiates a multi-agent self-learning harness that coordinates LLM and VLM agents to autonomously generate, simulate, evaluate, robot computation graphs in a loop with sampled task instances to increase a weighted combination of success rate and throughput until the combination reaches a plateau.  The resulting computation graph is then sent to an edge-based graph interpreter for repeated execution on the robot.    

%To prevent the generative agents from reward hacking from defining overly simplistic success criteria, an independent Evaluation Agent explicitly authors intermediate task checkpoints to enforce strict sub-task post-conditions.During simulated execution, the system provides precise state and contact feedback, such as verifying rigid-body interactions. Upon a checkpoint failure, the system localizes the error to a specific node and iteratively refines the graph until it achieves a plateau in internal simulation.
 
% GaP refines the policy graph through closed-loop internal simulation feedback. Given a task specification, a known distribution of initial states, and an initial policy graph, GaP first instantiates a parameterized internal physics simulation. To prevent the graph-generating LLM from optimizing toward overly simplistic checkpoints, GaP uses a separate LLM agent to explicitly define intermediate task checkpoints. These checkpoints establish post-conditions that must be satisfied for a task segment to be  successful. As the policy graph executes within the simulated environment, the internal simulation generates precise state and contact feedback, such as verifying whether a gripper has successfully established rigid-body contact with a target object. When a checkpoint is failed, a dedicated Coordinator Agent localizes the failure mode to a specific node and then iteratively refines the policy until the policy achieves a target success rate.
%\eric{discuss multi-agent to avoid cheating}


To evaluate GaP, we present 8 new open VA benchmarks. The first 6 -- based loosely on persistent commercial applications -- are (I-a,b) Fulfill Grocery Orders, (II-a,b) Pack Grocery Items, and (III-a,b) Make Popcorn, where a is in sim and b is in real. These 6 tasks use common kitchen items and one Franka robot arm from the LIBERO~\cite{liu2023libero} benchmark. The other two VA benchmarks -- based loosely on datacenter and industrial applications -- are (IV) Insert USB-C Cables, and (V) Wash Crates. Results suggest that GaP can achieve high success rates that significantly outperform baselines, including $\pi_{0.5}$~\cite{black2025pi_} and MolmoAct2~\cite{fang2026molmoact2}.
% contributon:

This paper makes the following contributions:
(1) the VA task class for robot learning with 8 open VA benchmarks; (2)  a computation graph structure for agentic robotics; % , a directed computation graph formulation that encapsulates atomic perception, planning, and control operations into interpretable, and modular robot policies.
(3) GaP, an implemented multi-agent harness that autonomously decomposes natural language specifications to collaboratively generate robot computation graphs;
(4) an open, evolving Modular Open Robot Skill Library (MORSL) with 51 initial skills;
(5) a self-learning approach to graph rehearsal and optimization using the Isaac-Lab physics simulator; and
(6) Experimental data comparing GaP with TAMP, model-free, and other baselines and simulations on the first 6 VA benchmarks, and performance results from GaP on the other 2 VA benchmarks.

% \eric{Note on GaP vs Multi-Agent vs Harness: GaP is a explicit hierarchical graph representation of robot policy, and using multiple agents and providing different agents different context (harness) is a method to generate GaP. I think these are two contributions. Also, I don't think harness is a term widely used in research community. I surveyed related work from ICLR/NeurIPS, they still uses multi-agent system / multi-agent collaboration / multi-agent workflow (optimization) in this context. 
% }

% \eric{Note on GaP vs coding agent: GaP is not a coding agent (although it was inspired by CAP-X initially). It uses LLM's generative and reasoning capability to generate a graph that specifies the different skills as nodes and their edge connectivity. }


% \eric{Note on the privileged information: CaP-X uses privileged information at policy inference stage. For example, originally the code needs to program a perception module itself to know the mask and 3D position, but with privileged information,  the code can just read the ground truth of the object location (that's why they call it privileged information). GaP uses internal simulation with Task Class (mostly object and scene geometry such as object mesh, object center pose belief space, etc) which makes it different from "privileged". }
